\documentclass[12pt]{article}

\usepackage[utf8]{inputenc}
\usepackage[T1]{fontenc}
\usepackage{lmodern}
\usepackage[a4paper,margin=1in]{geometry}
\usepackage{amsmath,amssymb,mathtools}
\usepackage{setspace,enumitem,microtype}
\usepackage{titlesec}

\setstretch{1.2}
\setlength{\parindent}{0pt}
\setlength{\parskip}{8pt}

\titleformat{\section}{\Large\bfseries}{}{0em}{}
\titleformat{\subsection}{\normalsize\bfseries}{}{0em}{}

\begin{document}

\begin{center}
    {\LARGE \textbf{Final Exam Review}}\\[4pt]
    \textit{Based on Alpha C. Chiang and Kevin Wainwright (4th Edition)}\\[4pt]
    Material: Functions, Linear Models, Differentiation, Optimization, Exponential \& Logarithmic Functions, Constrained Optimization, Envelope Theorem, Integration, Dynamics
\end{center}

\hrule
\vspace{10pt}

\section*{Functions and Graphs}

\textbf{Definitions}
\begin{itemize}
    \item A function assigns to each input \(x\) exactly one output \(y=f(x)\).
    \item Domain: all permissible \(x\)-values. Range: the set of resulting \(y\)-values.
\end{itemize}

\textbf{Common Function Types}
\[
\text{Linear: }y=a+bx,\qquad
\text{Quadratic: }y=ax^2+bx+c,\qquad
\text{Power: }y=ax^b.
\]
\[
\text{Exponential: }y=Ae^{bx},\qquad
\text{Logarithmic: }y=a+b\ln x.
\]

\textbf{Economic Uses}
\begin{itemize}
    \item Derivatives represent marginal changes (marginal cost, marginal utility).
    \item Elasticity measures percentage responsiveness.
\end{itemize}

\hrule
\vspace{10pt}

\section*{Linear Functions}

\textbf{General Form:}
\[
y = a + bx.
\]

\textbf{Slope:} \(b=\dfrac{\Delta y}{\Delta x}\).

\textbf{Applications}
\begin{itemize}
    \item Demand: \(Q_d=a-bP\)
    \item Supply: \(Q_s=c+dP\)
    \item Market equilibrium: \(Q_d=Q_s\).
\end{itemize}

\textbf{Elasticity}
\[
E = \frac{dy}{dx}\cdot \frac{x}{y}.
\]
Note: Elasticity varies along a linear function.

\hrule
\vspace{10pt}

\section*{Differentiation}

\textbf{Key Rules}
\[
\frac{d}{dx}x^n=nx^{n-1},\qquad
(uv)'=u'v+uv',\qquad
\left(\frac{u}{v}\right)'=\frac{u'v-uv'}{v^2}.
\]
\[
\text{Chain Rule: }(f(g(x)))'=f'(g(x))\cdot g'(x).
\]

\textbf{Economic Interpretation}
\begin{itemize}
    \item \(MC = C'(Q)\), \(MR = R'(Q)\).
    \item Profit maximization requires \(MR=MC\).
\end{itemize}

\hrule
\vspace{10pt}

\section*{Optimization of Single-Variable Functions}

To find extrema:
\[
f'(x)=0 \quad \text{(FOC)}.
\]

\textbf{Second-Order Test}
\begin{itemize}
    \item If \(f''(x)<0\): maximum.
    \item If \(f''(x)>0\): minimum.
\end{itemize}

\textbf{Inflection Point:} when \(f''(x)=0\) and concavity changes sign.

\textbf{Examples}
\begin{itemize}
    \item Profit maximization \(\pi(Q)=R(Q)-C(Q)\).
    \item Cost minimization with single variable.
\end{itemize}

\hrule
\vspace{10pt}

\section*{Exponential and Logarithmic Functions}

\textbf{Exponential Form}
\[
y=Ae^{bx}\quad\Rightarrow\quad y' = by.
\]
Proportional growth rate is constant:
\[
\frac{y'}{y}=b.
\]

\textbf{Logarithmic Function}
\[
y=a+b\ln x \quad\Rightarrow\quad y'=\frac{b}{x}.
\]

\textbf{Elasticity}
\[
E = \frac{dy}{dx}\cdot\frac{x}{y}.
\]
Log-linear forms have constant elasticity.

\textbf{Economic Applications}
\begin{itemize}
    \item Continuous compounding: \(A=Pe^{rt}\).
    \item Growth models: \(Y(t)=Y_0 e^{rt}\).
\end{itemize}

\hrule
\vspace{10pt}

\section*{Optimization of Multi-Variable Functions}

\textbf{First-Order Conditions}
\[
f_x(x,y)=0,\qquad f_y(x,y)=0.
\]

\textbf{Hessian Matrix}
\[
H = 
\begin{bmatrix}
f_{xx} & f_{xy}\\
f_{yx} & f_{yy}
\end{bmatrix}.
\]

\textbf{Classification}
\begin{itemize}
    \item If \(\det H>0\) and \(f_{xx}<0\): maximum.
    \item If \(\det H>0\) and \(f_{xx}>0\): minimum.
    \item If \(\det H<0\): saddle point.
\end{itemize}

\textbf{Economic Interpretations}
\begin{itemize}
    \item Diminishing marginal returns.
    \item Substitutability from cross-partials.
\end{itemize}

\hrule
\vspace{10pt}

\section*{Optimization with Equality Constraints}

Maximize \(f(x,y)\) subject to \(g(x,y)=b\).

\textbf{Lagrangian}
\[
\mathcal{L}(x,y,\lambda)=f(x,y)-\lambda(g(x,y)-b).
\]

\textbf{FOCs}
\[
f_x = \lambda g_x,\qquad f_y = \lambda g_y,\qquad g(x,y)=b.
\]

\textbf{Interpretation}
\(\lambda\) is the shadow price: the marginal value of relaxing the constraint.

\textbf{Bordered Hessian}
\[
H_B =
\begin{bmatrix}
0 & g_x & g_y\\
g_x & f_{xx} & f_{xy}\\
g_y & f_{xy} & f_{yy}
\end{bmatrix}.
\]
Sign of \(\det(H_B)\) determines max/min.

\hrule
\vspace{10pt}

\section*{Envelope Theorem and Duality}

\textbf{Envelope Theorem}
If \(v(b)\) is the maximum of \(f(x)\) subject to \(g(x)=b\), then
\[
\frac{dv}{db}=\lambda^\star.
\]
This simplifies comparative statics by avoiding differentiation of optimal choices.

\textbf{Duality}
\begin{itemize}
    \item Utility maximization vs cost minimization.
    \item Both yield the same marginal rate of substitution conditions.
    \item Shadow prices appear in both problems.
\end{itemize}

\hrule
\vspace{10pt}

\section*{Integration}

\textbf{Basic Rules}
\[
\int x^n dx = \frac{x^{n+1}}{n+1}+C,\qquad
\int e^{ax} dx = \frac{1}{a}e^{ax}+C,\qquad
\int \frac{1}{x}dx = \ln|x|+C.
\]

\textbf{Definite Integral}
\[
\int_a^b f(x)dx = F(b)-F(a).
\]

\textbf{Economic Uses}
\begin{itemize}
    \item Consumer surplus.
    \item Present value: \(\int_0^\infty e^{-rt}R(t)\,dt\).
    \item Capital accumulation solved by integration.
\end{itemize}

\hrule
\vspace{10pt}

\section*{Economic Dynamics}

\textbf{Differential Equations}
\[
\dot y + ay = b \quad \Rightarrow\quad y(t)=Ce^{-at}+\frac{b}{a}.
\]

\textbf{Growth Models}
\[
Y(t)=Y_0 e^{rt}.
\]

\textbf{Decay Models}
\[
K(t)=K_0 e^{-\delta t}.
\]

Dynamics describe continuous-time adjustment processes in economics (capital, population, debt).

\hrule
\vspace{10pt}

\section*{Summary Table}

\begin{tabular}{p{3cm} p{4cm} p{4.5cm}}
\textbf{Topic} & \textbf{Mathematical Tool} & \textbf{Economic Interpretation} \\
\hline
Slopes & Derivative & Marginal change \\
Concavity & Second derivative & Diminishing returns \\
Elasticity & Log derivative & Responsiveness \\
Multivariate optimization & Hessian matrix & Classification of behavior \\
Constraints & Lagrange multipliers & Shadow values \\
Comparative statics & Envelope theorem & Value-function sensitivity \\
Integration & Antiderivatives & Accumulation, surplus \\
Dynamics & Differential equations & Growth, decay, adjustment paths \\
\end{tabular}

\end{document}
